% =====================================================================
%  PHLOP -- updated Methodology / Results text
%  Reflects the changes made AFTER the original submission:
%    * train/val/test domain-generalization split (phlop/split_config.py)
%    * per-question difficulty labels      (phlop/question_answer.py,
%                                            phlop/advanced_physics_questions.py)
%    * difficulty-based fine-tuning ablation (notebooks/phlop_eval_common.py
%                                             -> FINE_TUNE_CONFIGS)
%    * static / moving camera modes
%    * physics-hint prompting toggle for fine-tuning
%    * new, *fine-tuned* model set: SmolVLM, Qwen2-VL, InternVL3
%  Requires: \usepackage{booktabs}  \usepackage{multirow}
% =====================================================================

% ---------------------------------------------------------------------
% WHAT CHANGED vs. the original submission (for the author; delete before
% camera-ready). The original paper described a single pool of 4,000 videos
% evaluated *zero-shot* with four prompting conditions over three off-the-shelf
% models (InternVideo2.5-8B, InternVL3-2B, Video-LLaMA3-7B). The revision adds:
%   1. An explicit, disjoint train/val/test split for domain generalization.
%   2. A per-question difficulty annotation (easy/medium/hard/very_hard).
%   3. A LoRA fine-tuning study over four difficulty curricula, with validation
%      on the held-out difficulties, under two prompting settings.
%   4. A new fine-tuned model set (SmolVLM, Qwen2-VL, InternVL3).
% ---------------------------------------------------------------------


% =====================================================================
\subsection{Dataset Splits for Domain Generalization}
% (Insert into Section III-A, after "Parameter Sampling".)
% =====================================================================

To evaluate \emph{generalization} rather than memorization, PHLOP is partitioned
into \textsc{train}, \textsc{val}, and \textsc{test} splits whose generative
factors are deliberately \emph{disjoint}. Rather than randomly partitioning a
single pool of scenes, each split draws from a different region of the scene
parameter space: distinct object shapes, materials, per-material physical
sub-components (e.g.\ different metal alloys), object-count ranges, collision
modes, floor textures, and camera profiles. Consequently, a model is tested on
material--shape--viewpoint combinations it never saw during training
(e.g.\ \texttt{rubber} and \texttt{block} objects, and an ``extreme'' camera
profile, appear \emph{only} in the test split). Table~\ref{tab:splits}
summarizes the configuration; the full specification is released in
\texttt{split\_config.py}.

\begin{table}[t]
\centering
\caption{Disjoint generative factors per split, enforcing
domain-generalization evaluation. Material sub-components (e.g.\ metal alloys)
are also disjoint across splits, so even shared material categories use
different physical-parameter regions.}
\label{tab:splits}
\small
\begin{tabular}{llll}
\toprule
Factor & Train & Val & Test \\
\midrule
Shapes        & ball, cube, cyl.   & ball, cube         & block, cyl., cube \\
Materials     & metal, wood, plastic & metal, wood, glass & glass, rubber \\
\#Objects     & 2--6               & 2--5               & 3--7 \\
Camera        & broad              & narrow             & extreme \\
Collisions    & coll./slide/stat.  & coll./slide        & coll./offset/slide \\
Floor         & checker/flat       & flat/gradient      & checker/gradient \\
\bottomrule
\end{tabular}
\end{table}

Every scene is additionally rendered under two \textbf{camera modes}---a
\texttt{static} fixed-viewpoint render and a \texttt{moving} (orbiting) render---
allowing the same physical event to be probed under different visual coverage.
All experiments reported here use the \texttt{static} mode.


% =====================================================================
\subsection{Question Difficulty Annotation}
% (Insert into Section III-B, after "Automatic Question Generation".)
% =====================================================================

Each automatically generated question is tagged with a \textbf{difficulty}
label---\textit{easy}, \textit{medium}, \textit{hard}, or \textit{very\_hard}---
reflecting the depth of reasoning required rather than surface form. The label is
assigned by the question generator from the reasoning type:

\begin{itemize}
  \item \textbf{Easy}: single-step perception and counting (object count, rolling
        detection, stopped-object count, collision presence).
  \item \textbf{Medium}: property comparison and attribute inference
        (e.g.\ highest-friction object, velocity/friction scaling).
  \item \textbf{Hard} / \textbf{very\_hard}: multi-step physical reasoning over
        continuous dynamics (momentum/energy conservation, multi-hop causation,
        temporal-event ordering).
\end{itemize}

These labels are the basis of the fine-tuning curriculum in
Section~\ref{sec:finetune}.


% =====================================================================
\subsection{Fine-tuning Protocol and Difficulty Ablation}
\label{sec:finetune}
% (Replaces / extends Section III-B-3 "Evaluation Protocol".)
% =====================================================================

Whereas our initial evaluation was restricted to zero-shot inference, we now
\textbf{LoRA-fine-tune} three open vision--language models---\textbf{SmolVLM},
\textbf{Qwen2-VL}, and \textbf{InternVL3}---on PHLOP and study how the
\emph{difficulty composition} of the training data affects physical reasoning.
We define four training curricula (Table~\ref{tab:ftconfig}). To measure
generalization \emph{across difficulty}, each model is validated on the
difficulties \emph{held out} of its training set; the \textit{full} curriculum
trains on all difficulties with no difficulty held out.

\begin{table}[t]
\centering
\caption{Fine-tuning curricula over question difficulty. ``Validation'' uses the
difficulties excluded from training (\texttt{FINE\_TUNE\_CONFIGS}).}
\label{tab:ftconfig}
\small
\begin{tabular}{lll}
\toprule
Curriculum & Train difficulties & Validation difficulties \\
\midrule
\textit{easy}        & easy               & medium, hard, very\_hard \\
\textit{easy\_medium} & easy, medium       & hard, very\_hard \\
\textit{hard}        & hard               & easy, medium, very\_hard \\
\textit{full}        & all                & --- (no holdout) \\
\bottomrule
\end{tabular}
\end{table}

Each curriculum is run under two prompting settings: \textbf{no physics hint}
(video + question only) and \textbf{with physics hint} (object physical
parameters appended to the prompt). Together with the non-fine-tuned
\emph{zero-shot} baseline, this yields five conditions per model and prompting
setting, evaluated on the same held-out \textsc{test} split
($25{,}436$ questions; $22{,}635$ excluding the degenerate \textit{unknown}
bucket). Results are reported in Section~\ref{sec:finetune-results} and
Table~\ref{tab:finetune-ablation}.
